\section*{2025 MIT Science Bowl High School Invitational}
\subsection*{Round 9}

\subsubsection*{TOSS-UP}
\textbf{1) CHEMISTRY}
Short Answer In iodine heptafluoride, what is the bond angle in degrees between adjacent equatorial fluorine atoms?
\\ANSWER: 72

\subsubsection*{BONUS}
\textbf{1) CHEMISTRY}
Multiple Choice If the average lengths of a carbon carbon single bond and carbon carbon double bond are 154 and 134 picometers, respectively, which of the following is closest to the length in picometers of the carbon carbon bond in acetylene?
\\W) 110
\\X) 120
\\Y) 130
\\Z) 140
\\ANSWER: X) 120

\subsubsection*{TOSS-UP}
\textbf{2) PHYSICS}
Short Answer A cylindrical resistor is created with the two circular faces of the cylinder as the two terminals. If its diameter is tripled and its length halved, by what factor is the resistance multiplied?
\\ANSWER: 1/18

\subsubsection*{BONUS}
\textbf{2) PHYSICS}
Short Answer A uniform sphere is released from rest at the top of a 30 degree incline with height h. When the sphere reaches the bottom of the incline, it continues sliding up a frictionless 45 degree incline. In terms of h, what is the maximum height the sphere reaches on the 45 degree incline?
\\ANSWER: 5h/7 (READ: FIVE OVER SEVEN H)

\newpage

\subsubsection*{TOSS-UP}
\textbf{3) MATH}
Short Answer The point (-3, 8) is rotated $90^{\circ}$ (read: 90 degrees) clockwise around the origin, then reflected across the line $y=x.$ What is the coordinates of the image of these transformations?
\\ANSWER: (3,8)

\subsubsection*{BONUS}
\textbf{3) MATH}
Short Answer Consider a 5 by 5 square lattice of points, such that the spaces between these points form a 4 by 4 grid of squares. How many unordered pairs of points in the lattice have a distance of at least 5 units?
\\ANSWER: 10

\subsubsection*{TOSS-UP}
\textbf{4) BIOLOGY}
Short Answer Tyrosinase (read: tie-RAH-sin-ase) is lacking in individuals with what disorder which is characterized by a lack of melanin production in keratinocytes (read: kerr-uh-TEE-noh-sights)?
\\ANSWER: ALBINISM

\subsubsection*{BONUS}
\textbf{4) BIOLOGY}
Multiple Choice Stimulated emission depletion microscopy is a novel microscopy technique which utilizes selective deactivation of neighboring fluorophores to narrow down their effective fluorescence spot and overcome the diffraction limit. This is similar to the function of what group of cells in the eye?
\\W) Horizontal cells
\\X) Rods
\\Y) Ganglion cells
\\Z) Bipolar cells
\\ANSWER: W) HORIZONTAL CELLS

\newpage

\textbf{5) EARTH AND SPACE}
\subsubsection*{TOSS-UP}
Multiple Choice Which of the following soil horizons is the most likely to contain high levels of kaolinite?
\\W) O
\\X) A
\\Y) E
\\Z) B
\\ANSWER: Z) B

\subsubsection*{BONUS}
\textbf{5) EARTH AND SPACE}
Short Answer Rank the following four greenhouse gases by increasing 100-year global warming potential: 1) Carbon dioxide; 2) Methane; 3) Water vapor; 4) Nitrous oxide.
\\ANSWER: 3, 1, 2, 4

\newpage

\subsubsection*{TOSS-UP}
\textbf{6) ENERGY}
Multiple Choice Researchers at the Broad Institute are studying biological factors underlying the epithelial to mesenchymal transition in metastatic mesoderm-derived epithelial cancerous tissue. Which of the following types of cancer were the researchers studying?
\\W) Ductal carcinoma
\\X) Basal cell carcinoma
\\Y) Colon adenocarcinoma
\\Z) Renal cell carcinoma
\\ANSWER: Z) RENAL CELL CARCINOMA

\subsubsection*{BONUS}
\textbf{6) ENERGY}
Multiple Choice The Collins Lab used high-throughput screening to identify semapimod as a potential new drug to target antibiotic-resistant bacteria. Unlike traditional antibiotics, semapimod succeeds at disrupting which of the following processes in antibiotic-resistant bacteria?
\\W) Cellular respiration
\\X) Protein biosynthesis
\\Y) DNA replication
\\Z) Membrane permeabilization
\\ANSWER: Z) MEMBRANE PERMEABILIZATION

\subsubsection*{TOSS-UP}
\textbf{7) PHYSICS}
Short Answer Jacob discovers a material whose relative permitivity and relative permeability are both equal to 5. Giving your answer to one significant figure in scientific notation, and in meters per second, what is the speed of light in this medium?
\\ANSWER: $6\times10^{7}$ (READ: SIX TIMES TEN TO THE SEVENTH)

\subsubsection*{BONUS}
\textbf{7) PHYSICS}
Short Answer A mass on a string rotates in a vertical circle such that the tension in the string is zero when the mass is at its highest point. What is the ratio between the mass's speed at the lowest and highest points of the circle, respectively?
\\ANSWER: $\sqrt{5}$ (READ: THE SQUARE ROOT OF FIVE)

\newpage

\subsubsection*{TOSS-UP}
\textbf{8) ENERGY}
Short Answer Researchers at CSAIL's Computation Structures Group are developing memory-efficient numerical computing systems. Under the two's complement scheme, what is the maximum value of a signed 8-bit integer?
\\ANSWER: 127

\subsubsection*{BONUS}
\textbf{8) ENERGY}
Short Answer The Kulik Group discovered an alkane monooxygenase that catalyzes the defluorination of 1-fluorooctane to produce octanal. How many electrons are transferred per mole of 1-fluorooctane by this reaction?
\\ANSWER: 2

\textbf{9) EARTH AND SPACE}
\subsubsection*{TOSS-UP}
Multiple Choice Mick is taking a rare walk outside and sees that the sun has been dulled into a dot with no halo. Which of the following clouds is most likely responsible for this?
\\W) Cirrus
\\X) Cirrostratus
\\Y) Altostratus
\\Z) Nimbostratus
\\ANSWER: Y) ALTOSTRATUS

\subsubsection*{BONUS}
\textbf{9) EARTH AND SPACE}
Short Answer Greenschist is green because it contains high levels of what phyllosilicate mineral that is also common in low-grade metamorphic rocks like phyllite?
\\ANSWER: CHLORITE

\newpage

\subsubsection*{TOSS-UP}
\textbf{10) BIOLOGY}
Short Answer Most strains of syphillis are resistant to 50S inhibiting antibiotics such as macrolides because they contain mutations in what enzyme found in the large ribosomal subunit that catalyzes the formation of peptide bonds?
\\ANSWER: PEPTIDYL TRANSFERASE

\subsubsection*{BONUS}
\textbf{10) BIOLOGY}
Multiple Choice A population initially in Hardy-Weinberg equillibrium is fixed at a locus, until one individual introduces a mutant allele. Given that the allele is highly beneficial, and inherited autosomal recessively, which of the following best describes how the frequency of the new mutant allele will change?
\\W) Increasing rapidly at first, then slowly until fixation
\\X) Increasing slowly at first, then rapidly until fixation
\\Y) Sinusoidally increasing and decreasing
\\Z) Increasing slowly at first, then decreasing until equillibrium
\\ANSWER: X) INCREASING SLOWLY AT FIRST, THEN RAPIDLY UNTIL FIXATION

\subsubsection*{TOSS-UP}
\textbf{11) MATH}
Short Answer Simmons Hall has 5500 windows across 300 dorm rooms. What is the maximum number n such that at least one dorm room is guaranteed to have at least n windows?
\\ANSWER: 19

\subsubsection*{BONUS}
\textbf{11) MATH}
Short Answer A function is concave up and increasing on a given interval. Identify all of the following options that are guaranteed to be an overestimate for the area under the curve: 1) Left Riemann sum; 2) Right Riemann sum; 3) Trapezoidal sum.
\\ANSWER: 2 AND 3

\newpage

\subsubsection*{TOSS-UP}
\textbf{12) CHEMISTRY}
Multiple Choice Which of the following is the major product from reacting 2-butyne with hydrogen gas over a lead-poisoned palladium catalyst?
\\W) n-Butane
\\X) But-1-ene
\\Y) Cis-but-2-ene
\\Z) Trans-but-2-ene
\\ANSWER: Y) CIS-BUT-2-ENE

\subsubsection*{BONUS}
\textbf{12) CHEMISTRY}
Short Answer Identify all of the following three values of the angular momentum quantum number that are allowed for an electron occupying a hydrogen atom's 3p orbital: 1) $L=-2$; 2) $L=0;$ 3) $L=2$
\\ANSWER: NONE

\subsubsection*{TOSS-UP}
\textbf{13) BIOLOGY}
Multiple Choice Larry is studying fragile X syndrome, which follows an X-linked dominant inheritance that gets worse every generation. Which of the following causes could explain why this is true?
\\W) Trinucleotide repeat expansion
\\X) Hypermethylation of X chromosome
\\Y) Loss of function inversion
\\Z) Gain of function insertion
\\ANSWER: W) TRINUCLEOTIDE REPEAT EXPANSION

\subsubsection*{BONUS}
\textbf{13) BIOLOGY}
Short Answer In the evolution of vertebrates, order the following three traits by the order in which they evolved: 1) Amniotic egg; 2) Mineralized skeleton; 3) Limbs with digits.
\\ANSWER: 2, 3, 1

\newpage

\textbf{14) CHEMISTRY}
Multiple Choice Which of the following cycloalkanes releases the most heat per carbon atom when combusted?
\\W) Cyclopropane
\\X) Cyclopentane
\\Y) Cycloheptane
\\Z) Cyclononane
\\ANSWER: W) CYCLOPROPANE

\subsubsection*{BONUS}
\textbf{14) CHEMISTRY}
Multiple Choice In the most stable chair conformation of cis 1-fluoro-3-tert-butyl cyclohexane, which of the following are the positions of the fluoro and tert-butyl substituents, respectively?
\\W) Equatorial, equatorial
\\X) Equatorial, axial
\\Y) Axial, equatorial
\\Z) Axial, axial
\\ANSWER: W) EQUATORIAL, EQUATORIAL

\subsubsection*{TOSS-UP}
\textbf{15) MATH}
Short Answer How many positive integer factors does 111,111 have?
\\ANSWER: 32

\subsubsection*{BONUS}
\textbf{15) MATH}
Short Answer An isosceles trapezoid has three congruent sides and an angle measuring 60 degrees. If one of its sides has length 4, what are all possible values for its perimeter?
\\ANSWER: 10, 20

\newpage

\subsubsection*{TOSS-UP}
\textbf{16) PHYSICS}
Multiple Choice A series RLC circuit performs damped oscillations analogous to those of a damped spring-mass oscillator. If the reciprocal of capacitance corresponds to the spring constant in a spring-mass system, the inductance is analogous to what quantity?
\\W) Mass
\\X) Spring constant
\\Y) Drag coefficient
\\Z) Q factor
\\ANSWER: W) MASS

\subsubsection*{BONUS}
\textbf{16) PHYSICS}
Short Answer Jack shines a flashlight on a piece of metal and is disappointed to find that he does not observe the photoelectric effect. Identify all of the following three reasons that could explain why: 1) The light's intensity is too low; 2) The light's wavelength is too low; 3) The light must be directed normal to the metal's surface.
\\ANSWER: NONE

\subsubsection*{TOSS-UP}
\textbf{17) EARTH AND SPACE}
Multiple Choice The end of a protostar's evolution along the Hayashi track is marked by what event?
\\W) Nuclear fusion of helium-3 begins
\\X) Formation of a radiative core
\\Y) Formation of bipolar outflows
\\Z) Kelvin-Helmholtz contraction ceases
\\ANSWER: X) FORMATION OF A RADIATIVE CORE

\subsubsection*{BONUS}
\textbf{17) EARTH AND SPACE}
Short Answer The radii of the Moon and the Earth both double at constant density, while their average separation distance remains unchanged. By what factor does the average height of the ocean's tides change?
\\ANSWER: 16

\newpage

\subsubsection*{TOSS-UP}
\textbf{18) ENERGY}
Multiple Choice Physicists at the Laboratory for Nuclear Science created an animated video to help visualize simulations of heavy atomic nuclei. Which of the following lengthscales is most appropriate for the atomic nuclei in their video?
\\W) Micrometer
\\X) Nanometer
\\Y) Picometer
\\Z) Femtometer
\\ANSWER: Z) FEMTOMETER

\subsubsection*{BONUS}
\textbf{18) ENERGY}
Short Answer Researchers from the Freedman Lab studied spin-phonon coupling in an organometallic complex. Identify all of the following three values that could be the spin of a phonon: 1) 0; 2) $1/2$; 3) 1.
\\ANSWER: 1 AND 3

\subsubsection*{TOSS-UP}
\textbf{19) CHEMISTRY}
Multiple Choice Which of the following elements has the greatest second ionization energy?
\\W) Fluorine
\\X) Neon
\\Y) Sodium
\\Z) Magnesium
\\ANSWER: Y) SODIUM

\subsubsection*{BONUS}
\textbf{19) CHEMISTRY}
Short Answer A solution containing zinc (II) is electrolyzed by an 8 ampere current for 12 hours. To the nearest tenth of a mole, at most how much metallic zinc is deposited?
\\ANSWER: 1.8

\newpage

\textbf{20) BIOLOGY}
\subsubsection*{TOSS-UP}
Multiple Choice A patient suffers a localized stroke that destroys their entire thalamus. Which of the following senses would be LEAST affected?
\\W) Vision
\\X) Audition
\\Y) Olfaction
\\Z) Gustation
\\ANSWER: Y) OLFACTION

\subsubsection*{BONUS}
\textbf{20) BIOLOGY}
Short Answer Hybridomas (read: high-bruh-DOH-muhs) are a type of synthetic tumor used to produce monoclonal antibodies. They are made by fusing immortal myeloma cells with what immune cell type?
\\ANSWER: B CELL (ACCEPT: PLASMA B CELL)

\subsubsection*{TOSS-UP}
\textbf{21) EARTH AND SPACE}
Multiple Choice Which of the following characteristics is most similar between Kuiper belt objects and scattered disk objects?
\\W) Orbital period
\\X) Perihelion distance
\\Y) Orbital eccentricity
\\Z) Orbital inclination
\\ANSWER: X) PERIHELION DISTANCE

\subsubsection*{BONUS}
\textbf{21) EARTH AND SPACE}
Multiple Choice At which of the following redshift values would you be most likely to spot a quasar?
\\W) 0.5
\\X) 2
\\Y) 6
\\Z) 20
\\ANSWER: X) 2

\newpage

\subsubsection*{TOSS-UP}
\textbf{22) PHYSICS}
Multiple Choice A spherical dust particle has a terminal velocity through air of v. In terms of v, which of the following is closest to the terminal velocity through air of a spherical dust particle made from the same material but with twice the radius?
\\W) v
\\X) 1.4v
\\Y) 2v
\\Z) 4v
\\ANSWER: X) 1.4v

\subsubsection*{BONUS}
\textbf{22) PHYSICS}
Short Answer A planet has a semimajor axis of $2.5\times10^{11}$ (read: two point five times ten to the eleventh) meters and a semiminor axis of $2.4\times10^{11}$ (read: two point four times ten to the eleventh) meters. In meters and to two significant figures, what is the minimum distance between the planet and its star?
\\ANSWER: $1.8\times10^{11}$ (READ: ONE POINT EIGHT TIMES TEN TO THE ELEVENTH)

\textbf{23) MATH}
\subsubsection*{TOSS-UP}
Multiple Choice How many solutions x are there from 0 to 27 for the equation $sin~x+cos~x=\sqrt{}3$ (read: the sine of x plus the cosine of x equals the square root of three)?
\\W) 0
\\X) 1
\\Y) 2
\\Z) 4
\\ANSWER: W) 0

\subsubsection*{BONUS}
\textbf{23) MATH}
Short Answer Mingwen chooses three distinct numbers from the set of integers from 1 to 5, inclusive. What is the expected value of the range of his set?
\\ANSWER: 3

\newpage

\subsubsection*{TOSS-UP}
\textbf{24) ENERGY}
Short Answer Researchers at the Jagoutz Group studied the serpentenization (read: serpent-ehn-eye-zay-tion) of olivine in Mars's crust. Serpentenization is evidence for the ancient presence of what gas in Mars's atmosphere?
\\ANSWER: WATER VAPOR (ACCEPT: WATER, $H_{2}O$, DO NOT ACCEPT: $O_{2}$)

\subsubsection*{BONUS}
\textbf{24) ENERGY}
Multiple Choice Researchers at CSAIL's Spoken Language Systems Group recently published a study inferring from the joint distribution of a passage of text. Which of the following best characterizes the entropy of the joint distribution when all random variables are mutually independent?
\\W) At a global minimum
\\X) At a global maximum
\\Y) At a saddle point
\\Z) At a non-critical point
\\ANSWER: X) AT A GLOBAL MAXIMUM

\subsubsection*{TOSS-UP}
\textbf{25) CHEMISTRY}
Multiple Choice Which of the following is the highest occupied molecular orbital in a semiconductor?
\\W) Valence band
\\X) Conduction band
\\Y) Band gap
\\Z) Fermi level
\\ANSWER: W) VALENCE BAND

\textbf{25) CHEMISTRY}
\subsubsection*{BONUS}
Multiple Choice Triphenylamine has a much greater bond angle than ammonia, trimethylamine, and nitrogen trifluoride. Which of the following best explains why?
\\W) The nitrogen atom in triphenylamine is sp2 hybridized
\\X) Phenyl groups create significant steric hindrance
\\Y) Phenyl groups have a low effective electronegativity
\\Z) Phenyl groups have a high effective electronegativity
\\ANSWER: W) THE NITROGEN ATOM IN TRIPHENYLAMINE IS SP2 HYBRIDIZED